% !TeX root = ../main.tex
\addchap*{Abstract}
Text-to-image (T2I) models produce convincing images from short prompts, but a
convincing image is no evidence that the model knows the depicted entity. This
thesis measures the factual knowledge of three T2I models, GPT Image 2,
FLUX.2-dev, and FLUX.2-klein-4B, about 298 real-world entities: 100 curated
entities from four domains, each annotated with an expected type and a
distinctive visual feature, and 198 randomly sampled Wikipedia entities in three
popularity tiers. Every entity is generated from the prompt ``An image of
[entity].'' A vision-language model, used as a judge, first describes each image without
knowing which entity it should show and lists short statements about what is
visible, called propositions; a second call checks these statements against the
introduction of the entity's Wikipedia article. On the 289 entities
available for all models, the models draw the correct coarse type for 73--79\%
of the curated entities, but the distinctive feature for only 63\% (GPT
Image 2), 46\% (FLUX.2-dev), and 29\% (FLUX.2-klein-4B). Of the extracted
propositions, 9.0\%, 9.7\%, and 17.7\% contradict the reference; for GPT
Image 2 this share rises from 3\% for the most popular to 23\% for the least
popular entities. The most severe errors come from reading the entity name as
something else, such as a portrait of a man for the British bus model Dennis
Dart. Unlike language models, image generators cannot abstain, and their
knowledge is limited to what can be drawn. Their outputs should therefore not
be used as factual depictions without verification.